% GENERATED, edit the French unit or the translation table
\mxnewpage\section*{EV3-T1 -- Mock exam 2}\label{EV3-T1}
\addcontentsline{toc}{section}{EV3-T1 -- Mock exam 2}
\textbf{Practice paper}. Programme studied in chapters 1 to 10.\par Duration: 2 h. No calculator. Marked out of 20. 1--3 April 2027.
\subsection*{Part 1 -- Fluency (6 points)}
Twelve questions, 0.5 point each. Only one correct answer; no justification required. An incorrect or missing answer scores 0, with no loss of points.
\begin{enumerate}
\item Calculate $6+2\times3$.\par A. $12$\quad B. $24$\quad C. $11$\quad D. $15$
\item Calculate $\frac12+\frac13$.\par A. $\frac{2}{5}$\quad B. $\frac{5}{3}$\quad C. $\frac{5}{6}$\quad D. $\frac{1}{6}$
\item Calculate $\dfrac{2^6}{2^{4}}$.\par A. $1$\quad B. $11$\quad C. $2$\quad D. $4$
\item A value of 50 increases by 20 \%. Final value?\par A. $10$\quad B. $40$\quad C. $70$\quad D. $60$
\item A value of 80 is obtained after a 20 \% decrease. Initial value?\par A. $64$\quad B. $96$\quad C. $100$\quad D. $60$
\item A 25 \% increase then a 20 \% decrease: overall coefficient?\par A. $\frac{21}{20}$\quad B. $\frac{9}{10}$\quad C. $1$\quad D. $\frac{4}{5}$
\item The solution of $2x+6=0$ is\par A. $12$\quad B. $-12$\quad C. $-3$\quad D. $3$
\item The slope between $(0;6)$ and $(2;12)$ is\par A. $-3$\quad B. $3$\quad C. $2$\quad D. $6$
\item The mean of $6, 8, 10$ is\par A. $9$\quad B. $8$\quad C. $24$\quad D. $6$
\item A sector of $90^\circ$ represents what fraction of the total?\par A. $\frac{1}{2}$\quad B. $\frac{1}{3}$\quad C. $\frac{1}{4}$\quad D. $\frac{3}{4}$
\item A probability is $\dfrac{2}{5}$. Its complement is\par A. $\frac{5}{2}$\quad B. $\frac{2}{5}$\quad C. $\frac{3}{5}$\quad D. $\frac{1}{5}$
\item For $f(x)=x^2+6$, $f(-2)$ is\par A. $10$\quad B. $8$\quad C. $24$\quad D. $2$
\end{enumerate}
\mxnewpage\subsection*{Part 2 -- Independent exercises (14 points)}
Every step must be justified. Exact results are expected unless stated otherwise.
\mxneed{12}
\subsection*{Exercise 1 -- A compound gain (5 points)}
A game consists of two independent trials, each with success probability $\dfrac12$. Each success earns 3 monetary units; taking part costs 2 units. Let $X$ be the net gain, that is, the total received minus the entry fee paid.
\begin{enumerate}
\item Establish the probability distribution of the net gain $X$. \hfill(1 pt)
\item Calculate its expectation and variance. \hfill(2 pts)
\item Determine the entry fee that would make the game fair. Explain what the standard deviation of the net gain would become. \hfill(2 pts)
\end{enumerate}
\mxneed{12}
\subsection*{Exercise 2 -- Peak of a model (5 points)}
A quantity is modelled for $t\geq0$ by the function $f(t)=(t+2)e^{-t}$. The time $t$ is expressed in a fixed unit. We look for its maximum value, then an approximation near the initial time.
\begin{enumerate}
\item Differentiate f. \hfill(1 pt)
\item Study the variation and give the maximum. \hfill(2 pts)
\item Write the tangent at 0 and approximate $f(0.1)$. \hfill(2 pts)
\end{enumerate}
\mxneed{12}
\subsection*{Exercise 3 -- The unit circle (4 points)}
The plane is equipped with an orthonormal frame. On the unit circle with centre $O$, the points $A$ and $B$ are the images of the real numbers $\dfrac\pi6$ and $\dfrac{5\pi}6$ respectively.
\begin{enumerate}
\item Give their coordinates. \hfill(1 pt)
\item Calculate AB. \hfill(1 pt)
\item Give the equation of their perpendicular bisector and justify geometrically. \hfill(2 pts)
\end{enumerate}
